% E4: Ethics Statement
% Addresses ethical considerations specific to automated LLM routing.
% Complements E3 (Broader Impact) with concrete ethical guidelines.

\subsection{Ethics Statement}
\label{appendix:ethics}

Automated LLM routing raises ethical questions beyond the societal-level
concerns discussed in Appendix~\ref{appendix:broader_impact}.  This section
addresses considerations that practitioners should evaluate before deploying
banditGPT or similar routing systems.

\paragraph{User transparency and informed consent.}
banditGPT routes user queries to different LLM backends without requiring
user input on the routing decision.  In most deployments, users interact
with a single API endpoint and may not be aware that their prompt was served
by a cheaper or less capable model.  This raises a transparency concern:
users who pay for a premium service may expect a frontier model for every
query.

\emph{Recommendation:}  Deployments should disclose that adaptive routing
is in use, either through documentation, terms of service, or per-response
metadata indicating which model served the query.  The library's
\texttt{explain\_decision} API returns the selected model and routing
rationale, enabling operators to surface this information to users.
Where full disclosure is impractical, operators should at minimum guarantee
a quality floor ($\lambda{=}0$ or \texttt{quality\_floor} constraint) that
ensures no user receives systematically worse service than a static
baseline.

\paragraph{Data privacy and prompt handling.}
The routing pipeline processes user prompts through a sentence embedding
model (\texttt{all-MiniLM-L6-v2}) to compute context features.  While the
PCA-reduced embeddings are low-dimensional (32 dimensions) and not trivially
invertible to recover the original prompt, they do encode semantic content.
The system also stores prompt--reward pairs during online learning, which
may contain sensitive user information.

\emph{Recommendation:}  Operators should apply the same data handling,
retention, and privacy policies to the routing pipeline as to the LLM
inference pipeline itself.  The embedding computation runs locally (no
external API calls), and the \texttt{ContextStore} can be configured with
retention limits.  For privacy-sensitive deployments, operators can
pre-compute embeddings in a secure environment and discard raw prompts
before routing, or use differential privacy techniques on the stored
context vectors.

\paragraph{Fairness across user populations.}
If prompt characteristics correlate with user demographics (e.g., language,
domain expertise, cultural context), the router may learn to systematically
route certain user groups to cheaper models.  This is not an explicit design
choice but an emergent property of optimizing for reward under cost
constraints: if the judge rates certain prompt types as ``easy'' (high tie
rate between models), the router learns to route them to the cheaper model,
even if those users would benefit from the frontier model.

\emph{Recommendation:}  Before production deployment, audit the router's
model selection patterns stratified by relevant demographic or domain
categories.  The calibration analysis
(Appendix~\ref{appendix:calibration}) provides a template for checking
prediction bias per model.  Operators should extend this to check for
disparate routing rates across user segments.  The $\gamma$-mixing floor
in Corralling ensures no model is completely abandoned, but it does not
guarantee equitable routing rates across user groups --- explicit fairness
constraints may be needed for high-stakes applications.

\paragraph{Accountability for routing decisions.}
When a routing decision leads to a poor outcome (e.g., an incorrect or
harmful response from the cheaper model), accountability is diffused across
the routing system, the selected model, and the reward signal that trained
the router.  Unlike static model selection, where the operator explicitly
chooses which model to deploy, adaptive routing delegates this decision to
an algorithm.

\emph{Recommendation:}  Operators should maintain routing logs that record
the selected model, context features, and confidence score for each
request.  The \texttt{explain\_decision} API provides this information.
For high-stakes domains (medical, legal, financial), consider using the
router in an advisory capacity (suggesting a model) rather than making
fully autonomous routing decisions, or enforce \texttt{quality\_floor}
constraints that prevent routing to models below a validated performance
threshold.

\paragraph{Ethics of automated quality judgments.}
The routing system learns from an LLM judge (GPT-4o) that produces the
reward signal.  This creates an implicit ethical dependency: the router
optimizes for the judge's notion of quality, which may not align with all
users' values.  For example, the MT-Bench-style judge prompt emphasizes
helpfulness, accuracy, and detail, but does not explicitly evaluate safety,
cultural sensitivity, or accessibility.

\emph{Recommendation:}  Practitioners should evaluate whether the judge's
quality criteria match their deployment's ethical requirements.  The
banditGPT framework is reward-agnostic
(Appendix~\ref{appendix:reward_signal}): operators can replace or augment
the LLM judge with human evaluations, safety classifiers, or
domain-specific rubrics.  For deployments where safety is paramount,
consider a composite reward that includes a safety score alongside the
quality judgment.

\paragraph{Dataset ethics and provenance.}
The LMSYS evaluation data consists of real user conversations from the
Chatbot Arena platform~\cite{zheng2023judging}, which is a public
benchmark where users voluntarily submit prompts.  The 80K warmup battles
from RouteLLM are derived from the same platform.  We do not collect any
new user data; all experiments use existing public datasets.  However,
these datasets may contain prompts with personal information, offensive
content, or culturally sensitive material that was submitted by Chatbot
Arena users.

\emph{Recommendation:}  When using banditGPT with proprietary or
user-generated data, operators should follow their organization's IRB or
ethics review procedures.  The library does not require access to prompt
content during inference (only embeddings), which provides a degree of
separation between routing decisions and raw user data.

\paragraph{Environmental cost of experimentation.}
All experiments in this paper were conducted using pre-computed LLM
responses and judge verdicts --- no new LLM inference calls were made
during the bandit experiments themselves.  The total compute for all
appendix experiments (1,400+ ablation runs, 360 PCA-vs-K runs, learning
curves, statistical tests) was under 5 minutes on a single Apple M2
laptop.  The embedding computation (SentenceTransformer) adds approximately
4 seconds per experiment.  We consider the environmental footprint of the
experimental evaluation to be minimal.
